\section{Experimental Setup}
\label{sec:setup}

\subsection{Timed stages}
We time four cumulative stages (defined and motivated in Section~\ref{sec:bg-stages}) on the same batch of
collocation points:
\begin{enumerate}
  \item \textbf{forward}: evaluate $\Th(x)$, no derivatives -- a plain prediction;
  \item \textbf{first\_grad}: the forward pass plus $\nabla_x\Th$;
  \item \textbf{second\_grad}: the forward pass, $\nabla_x\Th$ and $\partial^2\Th/\partial z^2$ -- everything
  the PDE residual $r_\theta$ needs;
  \item \textbf{param\_grad}: the second\_grad graph plus one further backward pass of the PDE-residual loss
  to the parameters, $\partial\Lloss_{\text{PDE}}/\partial\theta$ (Section~\ref{sec:bg-param}) -- approximately
  the PDE term's share of one optimiser step.
\end{enumerate}
The stages are cumulative, so a ratio such as the \emph{second-over-forward ratio} is the cost of everything
needed for the PDE residual relative to a plain evaluation, not the cost of the second derivative alone.
param\_grad covers only the PDE term of $\Lloss$; the boundary, initial-condition and data terms add further
backward passes that we do not separately time, so a full training step costs somewhat more than param\_grad.

\subsection{Size-matched classical controls}
\label{sec:setup-controls}
For each hybrid configuration we also time a purely classical network in which the VQC is replaced by a
$\tanh$ layer of the same width $n$, giving a network with exactly the hybrid's classical parameter count
(\texttt{ctrl\_n2}, \texttt{ctrl\_n4}, \texttt{ctrl\_n6}; \emph{deep} shares its control with \emph{baseline},
since both have $n=4$). Comparing a hybrid to its matched control, rather than only to the full classical
baseline, separates the cost of \emph{having fewer parameters} from the cost of \emph{simulating the quantum
circuit}: the matched control has the hybrid's parameter count but no circuit to simulate.

\subsection{Protocol}
Each stage is timed over 5--20 repetitions per round (adaptively chosen so a round takes about a fixed time
budget) after 3 warm-up calls, on a batch of 2540 collocation points sampled uniformly in $[0,2]\times[0,5]$,
matching the training batch size; batch size is itself swept from 100 to 5000 in a separate experiment
(Section~\ref{sec:res-batch}). Before the first timed section, a 120\,s settle phase of untimed work on
the drift-check cases brings CPU clocks and thread scheduling to a steady state; this was added after an
earlier run showed the hybrid configurations running up to 9\% faster at the end than at the start, consistent
with CPU warm-up rather than background load. Cases (configuration $\times$ seed) are interleaved in a shuffled
order within each round, so that slow drift of the machine affects all cases alike rather than biasing whichever
ran last. Network weights, batch data and circuit parameters are redrawn for each of 3 independent seeds
(0, 1, 2); we report the mean over seeds of the per-seed mean, and the sample standard deviation across seeds.
Backends are \texttt{default.qubit} (vectorised over the batch) and \texttt{lightning.qubit}, compared on the
circuit alone (Section~\ref{sec:res-backend}).

A drift check re-times a subset of the \emph{main} cases (\emph{classical}, \emph{baseline}, \emph{wide}) at the
end of the run and compares them against the same cases at the start. Judged against all cells, this flagged
deviations above 10\% on every attempt, but the failures were driven almost entirely by the classical model's
sub-millisecond stages, where a handful of scheduling-jitter outliers move the mean by 10--30\% while the
minimum barely changes. The check now uses the median ratio of only the cells that take at least 5\,ms in
\emph{main}; the faster, jitter-dominated cells are still reported but excluded from the verdict. Under this
rule, the sixth independent run (\texttt{win\_py314\_run6}) -- taken with the machine on a high-performance
power plan, on AC power, with background applications closed, and after the settle phase above -- passes with
a largest deviation of 4.2\% (\emph{wide} \texttt{second\_grad}), the smallest of six independent runs, and its
numbers are the ones reported in this paper.
